% Generated draft: requires mathematical and source review.
\documentclass[11pt]{article}
\usepackage[margin=25mm]{geometry}
\usepackage{fontspec}
\setmainfont{Latin Modern Roman}
\usepackage{amsmath,amssymb,mathtools,unicode-math}
\setmathfont{Latin Modern Math}
\usepackage{enumitem,needspace,xurl,hyperref}
\setlength{\parindent}{0pt}
\setlength{\parskip}{5pt}
\setlength{\emergencystretch}{3em}
\newcommand{\N}{\mathbb{N}}
\newcommand{\Z}{\mathbb{Z}}
\newcommand{\Q}{\mathbb{Q}}
\newcommand{\R}{\mathbb{R}}
\newcommand{\C}{\mathbb{C}}
\newcommand{\F}{\mathbb{F}}
\newcommand{\A}{\mathbb{A}}
\newcommand{\PP}{\mathbb{P}}
\newcommand{\E}{\mathbb{E}}
\newcommand{\eps}{\varepsilon}
\newcommand{\dd}{\,\mathrm{d}}
\begin{document}
\section*{UnsolvedMath ID: 800336}
\textbf{Source problem code:} catalogue entry
\subsection{Definitions and mathematical statement}
A matching of a graph is a set of edges no two of which share a vertex; it is perfect if every vertex belongs to one of those edges. Write \(H_{a,b,c}=(V,E)\) for the intended \(a,b,c\) honeycomb graph. The supplied record does not define this graph, the admissible parameters, or its embedding and boundary conditions.

For a specified family \(\mathcal{M}\) of matchings of \(H_{a,b,c}\) and a specified probability law \(\mathbb{P}\) on that family, define the edge probability by
\[
p_e=\mathbb{P}\{M\in\mathcal{M}:e\in M\},\qquad e\in E.
\]
The source does not specify whether these are perfect matchings or whether the probability law is uniform.

\textbf{Requested determination (Propp, Problem 7, 1998).} Find the ``vertical moments of inertia'' of the mass on edges arising from these edge probabilities, for the intended \(a,b,c\) honeycomb graphs.

The phrase defining the requested moments requires further specification. For example, if each edge probability is a point mass at a specified location \(r_e=(x_e,y_e)\), then the moment about a specified vertical line \(L=\{(x,y):x=x_0\}\) would be
\[
I_L=\sum_{e\in E}p_e\,\operatorname{dist}(r_e,L)^2
=\sum_{e\in E}p_e(x_e-x_0)^2.
\]
If instead mass is distributed along edges, its distribution must be specified to define the corresponding integral. A second moment of the vertical coordinates would involve \((y_e-y_0)^2\) and is a different possible interpretation. None of these conventions is selected by the supplied excerpt. Thus the graph, matching law, mass placement, geometric normalization, and intended meaning of ``vertical'' remain unresolved; the displayed formula is conditional and is not asserted to be the source's intended formulation.
\subsection{Short English statement}
Propp's Problem 7 asks for the vertical moments of inertia of edge-probability masses in random matchings of an \(a,b,c\) honeycomb graph. The supplied excerpt leaves the graph, probability model, mass placement, and moment convention unspecified, so it does not determine a unique mathematical quantity. The registry labels the problem open; that label is supplied metadata, not an independently verified status assessment.
\subsection{Sources}
\begin{itemize}
\item Supplied local UnsolvedMath registry record, permanent ID \(800336\), exact source code \texttt{catalogue entry}, titled ``Enumeration of Matchings --- Problem 7,'' in the AMR Open Problem Lists collection. Upstream revision: \texttt{d371ee917ac9647fbb06cd6fc529deec6f3dda82}. The precise local registry filename is not supplied.
\item Propp, \emph{Twenty Open Problems in Enumeration of Matchings} (1998), Problem 7 (\#7), as attributed in the supplied record: \url{https://arxiv.org/abs/math/9801060}. The background records access on 2026-10-16 and extraction by \texttt{arxiv\_latex}; these are supplied provenance statements, not a claim of independent source access or verification.
\item The supplied registry labels the problem \texttt{open} and records an update on 2026-10-16. No independent literature assessment or later resolution is supplied.
\end{itemize}
\end{document}
